% ============================================================================
% CHAPTER 1: INTRODUCTION
% ============================================================================

\chapter{Introduction}
\label{ch:introduction}

\section{Motivation}

The proliferation of Unmanned Aerial Vehicles (UAVs) in safety-critical applications---from urban air mobility to medical delivery, infrastructure inspection to search and rescue---demands control systems that can provide \textbf{mathematically verifiable safety guarantees}. Traditional approaches to drone control rely on heuristic safety mechanisms or ``black-box'' AI systems that cannot provide formal assurances about their behavior.

This document presents a fundamentally different approach: a \textbf{Contract-Based Adaptive UAV Control System} that leverages formal methods from computer science to provide provable safety properties while maintaining the adaptability required for real-world operation.

\begin{figure}[htbp]
    \centering
    \begin{tikzpicture}[
        scale=0.9,
        approach/.style={
            draw,
            rounded corners,
            minimum width=4cm,
            minimum height=2.5cm,
            align=center,
            font=\small
        }
    ]
        % Traditional approach
        \node[approach, fill=red!20] (trad) at (0,0) {
            \textbf{Traditional}\\[3pt]
            Heuristic safety\\
            No formal guarantees\\
            Reactive adaptation
        };

        % AI approach
        \node[approach, fill=orange!20] (ai) at (5.5,0) {
            \textbf{AI-Based}\\[3pt]
            Learned behavior\\
            Black-box decisions\\
            Opaque reasoning
        };

        % Our approach
        \node[approach, fill=green!20, line width=2pt, draw=contractpurple] (ours) at (11,0) {
            \textbf{Contract-Based}\\[3pt]
            Formal verification\\
            Provable guarantees\\
            Predictive adaptation
        };

        % Labels
        \node[above=0.5cm of trad, font=\footnotesize\itshape] {PX4, ArduPilot};
        \node[above=0.5cm of ai, font=\footnotesize\itshape] {Skydio, DJI};
        \node[above=0.5cm of ours, font=\footnotesize\itshape\color{contractpurple}] {This Work};

        % Progress arrow
        \draw[-{Stealth}, very thick, gray] (-2,-1.8) -- (13,-1.8)
            node[midway, below, font=\small] {Increasing Safety Assurance};
    \end{tikzpicture}
    \caption{Evolution of UAV control approaches: from heuristic to formal methods}
    \label{fig:approaches}
\end{figure}

\section{The Contract-Based Approach}

At the heart of our system lies the concept of \textbf{Assume-Guarantee (A/G) contracts}. Every component in the system operates under a formal contract that specifies:

\begin{definition}[Assume-Guarantee Contract]
An Assume-Guarantee contract $\mathcal{C} = (\mathcal{A}, \mathcal{G})$ consists of:
\begin{itemize}
    \item \textbf{Assumptions} $\mathcal{A}$: Conditions the environment must satisfy
    \item \textbf{Guarantees} $\mathcal{G}$: Properties the component will ensure
\end{itemize}
The contract semantics are: if $\mathcal{A}$ holds, then $\mathcal{G}$ is guaranteed.
\end{definition}

\importantbox{
\textbf{Key Insight:} Contracts enable \textit{compositional reasoning}---we can verify system-level properties by analyzing component contracts, without needing to understand internal implementations.
}

\section{System Overview}

The complete control system follows a hierarchical pipeline architecture:

\begin{figure}[htbp]
    \centering
    \begin{tikzpicture}[
        node distance=1.2cm and 0.8cm,
        block/.style={
            draw,
            rounded corners=4pt,
            minimum width=2.2cm,
            minimum height=1.4cm,
            align=center,
            font=\small,
            drop shadow={shadow xshift=1pt, shadow yshift=-1pt}
        },
        contract/.style={
            draw,
            dashed,
            rounded corners=2pt,
            fill=contractpurple!10,
            minimum height=0.6cm,
            font=\scriptsize,
            inner sep=3pt
        },
        arrow/.style={-{Stealth}, thick},
        dataarrow/.style={-{Stealth}, thick, blue!60}
    ]
        % Main pipeline blocks
        \node[block, fill=sensorblue!30] (sensors) {
            \textbf{Sensors}\\
            \scriptsize GPS, IMU
        };

        \node[block, fill=ekfgreen!30, right=of sensors] (ekf) {
            \textbf{EKF}\\
            \scriptsize State Est.
        };

        \node[block, fill=contractpurple!30, right=of ekf] (planner) {
            \textbf{Horizon}\\
            \textbf{Planner}\\
            \scriptsize Pacti
        };

        \node[block, fill=supervisororange!30, right=of planner] (supervisor) {
            \textbf{Flight Mode}\\
            \textbf{Supervisor}
        };

        \node[block, fill=pidblue!30, below right=1cm and -0.5cm of supervisor] (pid) {
            \textbf{PID}\\
            \scriptsize Efficient
        };

        \node[block, fill=hinfgreen!30, below left=1cm and -0.5cm of supervisor] (hinf) {
            \textbf{H-$\infty$}\\
            \scriptsize Robust
        };

        \node[block, fill=cbforange!30, below=2.5cm of supervisor] (cbf) {
            \textbf{CBF Filter}\\
            \scriptsize Safety
        };

        \node[block, fill=gray!30, below=of cbf] (actuators) {
            \textbf{Actuators}\\
            \scriptsize Motors
        };

        % Contract labels
        \node[contract, above=0.3cm of sensors] {$\mathcal{C}_{\text{sensor}}$};
        \node[contract, above=0.3cm of ekf] {$\mathcal{C}_{\text{estimator}}$};
        \node[contract, above=0.3cm of planner] {$\mathcal{C}_{\text{planner}}$};
        \node[contract, above=0.3cm of supervisor] {$\mathcal{C}_{\text{supervisor}}$};
        \node[contract, right=0.2cm of pid] {$\mathcal{C}_{\text{PID}}$};
        \node[contract, left=0.2cm of hinf] {$\mathcal{C}_{\text{H}\infty}$};
        \node[contract, right=0.3cm of cbf] {$\mathcal{C}_{\text{safety}}$};

        % Main flow arrows
        \draw[arrow] (sensors) -- (ekf);
        \draw[arrow] (ekf) -- (planner);
        \draw[arrow] (planner) -- (supervisor);

        % Controller selection
        \draw[arrow] (supervisor) -- (pid);
        \draw[arrow] (supervisor) -- (hinf);
        \draw[arrow] (pid) -- (cbf);
        \draw[arrow] (hinf) -- (cbf);
        \draw[arrow] (cbf) -- (actuators);

        % Feedback
        \draw[arrow, dashed, gray]
            (actuators.west) -- ++(-5,0) |- (sensors.west);

        % Selection indicator
        \node[font=\scriptsize, text=contractpurple] at ($(supervisor)!0.5!(pid) + (0.8,0)$) {select};
        \node[font=\scriptsize, text=contractpurple] at ($(supervisor)!0.5!(hinf) + (-0.8,0)$) {or};

    \end{tikzpicture}
    \caption{Complete system architecture with contract interfaces at each component boundary}
    \label{fig:system_architecture}
\end{figure}

\section{Key Contributions}

This work makes the following contributions to the field of autonomous UAV control:

\begin{enumerate}[label=\textbf{C\arabic*:}, leftmargin=*]
    \item \textbf{Hierarchical Contract Framework:} A complete specification of A/G contracts for all system components, enabling compositional verification from sensors to actuators.

    \item \textbf{Predictive Horizon Planning:} A novel use of contract composition over a planning horizon to anticipate safety violations \textit{before} they occur, enabling preemptive adaptation.

    \item \textbf{Principled Controller Switching:} Mathematically justified transitions between PID and H-infinity controllers based on contract satisfaction, not heuristics.

    \item \textbf{Multi-Layer Safety Architecture:} Defense-in-depth with contracts at the planning level, mode switching at the supervisor level, and CBF enforcement at the control level.

    \item \textbf{Pre-flight Verification:} The ability to mathematically prove mission feasibility before takeoff through contract composition.
\end{enumerate}

\section{Document Organization}

The remainder of this document is organized as follows:

\begin{description}[leftmargin=1.5cm, style=nextline]
    \item[Chapter 2:] \textbf{System Architecture} --- Detailed description of the control pipeline, data flow, and coordinate systems.

    \item[Chapter 3:] \textbf{Contract Framework} --- Comprehensive specification of A/G contracts, the Pacti library, and contract composition.

    \item[Chapter 4:] \textbf{State Estimation} --- The contract-aware Extended Kalman Filter with sensor fusion strategies.

    \item[Chapter 5:] \textbf{Control Systems} --- PID and H-infinity controllers, their contracts, and switching logic.

    \item[Chapter 6:] \textbf{Horizon Planning} --- Predictive planning using contract cascade over planning horizons.

    \item[Chapter 7:] \textbf{Safety Layer} --- Control Barrier Functions and runtime safety enforcement.

    \item[Chapter 8:] \textbf{System Integration} --- Complete adaptive control system implementation.

    \item[Chapter 9:] \textbf{Experimental Results} --- Multi-waypoint mission demonstrations and performance analysis.

    \item[Chapter 10:] \textbf{Market Analysis} --- Comparison with commercial and open-source alternatives.

    \item[Chapter 11:] \textbf{Conclusion} --- Summary and future directions.
\end{description}

\section{Target Applications}

The contract-based approach is particularly suited for safety-critical UAV applications:

\begin{table}[htbp]
    \centering
    \caption{Target applications and their safety requirements}
    \label{tab:applications}
    \begin{tabular}{@{}lll@{}}
        \toprule
        \textbf{Application} & \textbf{Safety Requirement} & \textbf{Contract Benefit} \\
        \midrule
        Urban Air Mobility & Passenger safety certification & Formal verification \\
        Medical Delivery & Guaranteed delivery reliability & Pre-flight verification \\
        Infrastructure Inspection & Predictable near-structure behavior & Bounded tracking error \\
        Search \& Rescue & Operation in degraded conditions & Graceful degradation \\
        Military/Defense & Certified autonomous behavior & Transparent decisions \\
        \bottomrule
    \end{tabular}
\end{table}
